% Source: https://artofproblemsolving.com/community/c3h2223910_local_maxima_and_minima_on_straight_line
% Stable identity: aops-topic-2223910-q01
\textbf{(a)} Let $a$ be any real number and take $I=(a-1,a+1)$. For every $x$ in $I$ we have $f(x)=1=f(a)$, so both $f(x)\le f(a)$ and $f(x)\ge f(a)$ hold. Hence $f$ has a local maximum \textbf{and} a local minimum at every point. (So the intuition in the question is right: with these standard definitions, which use $\le$ and $\ge$, every point of a horizontal line is both.)

\textbf{(b)} Any open interval $I$ containing $a$ contains some $x\ne a$, and for it $f(x)=f(a)$. So $f(x)<f(a)$ fails and $f(x)>f(a)$ fails: no point is a strict local maximum or a strict local minimum.

This is why some books speak of ``strict'' extrema, or of turning points, when they want to exclude flat stretches of a graph.
